% TOP_PROBLEM: {"id": 30007169, "title": "Dimension of a real self-similar attractor without exact overlaps", "category_id": 11, "category": {"id": 11, "name": "dynamical_systems", "display_name": "Dynamical Systems", "description": "Problems about long-term behavior of deterministic systems, Hamiltonian dynamics, and periodic orbits.", "slug": "dynamical-systems", "order_index": 11, "created_at": "2026-07-31T15:26:25.671Z"}, "status": "open"}
% FIRST_POSTED: 2026-10-07
% ATTEMPT_STATUS: solved
% ENTRY_KIND: human_contribution
\documentclass[11pt]{article}
\usepackage[margin=25mm]{geometry}
\usepackage{fontspec}
\setmainfont{Latin Modern Roman}
\usepackage{amsmath,amssymb,mathtools}
\usepackage{unicode-math}
\setmathfont{Latin Modern Math}
\usepackage{enumitem}
\usepackage{xurl,hyperref,bookmark}
\hypersetup{colorlinks=true,urlcolor=blue}
\setcounter{secnumdepth}{0}
\setlength{\parindent}{0pt}
\setlength{\parskip}{5pt}
\setlength{\emergencystretch}{3em}
\begin{document}
\section*{Human Contribution: Samuel Kittle and Constantin Kogler}

% Problem #2713 / local ID 30007169
% Insert immediately BEFORE \end{document} in the existing OpenAI
% contribution's LaTeX source. This is an \input-ready fragment:
% do not add a second preamble or \begin{document}.
% The host document already loads amsmath, amssymb, xurl, and hyperref.

\subsection*{Related independent contribution: Kittle--Kogler (2026)}

\paragraph{Paper and relation to this problem.}
Samuel Kittle and Constantin Kogler, \emph{The Exact Overlaps
Conjecture for Self-Similar Measures on the Real Line},
\href{https://arxiv.org/abs/2610.10511}{arXiv:2610.10511v1}
(posted October 7, 2026), give an independent proof of the
one-dimensional entropy-rate dimension formula. Their Corollary~1.2
settles exactly the attractor question recorded here: for a finite family
$g_i(x)=r_i x+b_i$ with $0<|r_i|<1$ and no exact overlaps, if $K$ is the
attractor and $s\geq 0$ satisfies $\sum_i|r_i|^s=1$, then
\[
  \boxed{\dim_H K=\min\{1,s\}.}
\]
This imposes no exponential separation or algebraicity hypothesis and
allows unequal and orientation-reversing contraction ratios.

\paragraph{Stronger measure theorem (Theorem 1.1).}
For a finitely supported probability law $\theta=\sum_i p_i\delta_{g_i}$
on contracting real similarities, let $\nu$ denote its stationary
self-similar measure. Write
\[
 h_\theta=\lim_{n\to\infty}\frac{H(\theta^{*n})}{n},
 \qquad
 \chi_\theta=-\sum_i p_i\log|r_i|>0,
\]
where $H(\theta^{*n})$ counts the entropy of the \emph{distinct composed
affine maps}, not merely symbolic words. Kittle--Kogler prove
\[
 \boxed{\dim_H\nu=\min\left\{1,\frac{h_\theta}{\chi_\theta}\right\}}
\]
with exact overlaps permitted. This has the same conclusion and scope
as the main theorem of OpenAI's Result~148. If there are no exact
overlaps, $h_\theta=H(p)$; choosing $p_i=|r_i|^s$ gives
$H(p)=s\chi_\theta$, yielding the lower bound for $\dim_H K$, while
the standard cylinder cover gives the matching upper bound.
Thus neither contribution leaves a case of the recorded one-dimensional
attractor problem unresolved by the other.

\paragraph{Additional result: entropy loss under addition (Theorem 1.3).}
A distinct contribution of Kittle--Kogler is a quantitative inequality
for arbitrary independent, finitely supported real random variables.
For a compactly supported probability law $\rho$, put
\[
 V_\rho(r)=\frac{4}{r^3}\int_{\mathbb R}
   \inf_{c\in\mathbb R}\int_{[a,a+r)}(x-c)^2\,d\rho(x)\,da,
 \qquad
 \mathcal E(\rho)=\int_0^\infty V_\rho(r)\,\frac{dr}{r}.
\]
Here $\mathcal E$ is the \emph{variance energy}, which accumulates
normalized local variance across logarithmic scales. For independent
finitely supported $X_1,\ldots,X_m$, their Theorem~1.3 gives
\[
 \boxed{
 \sum_{j=1}^m H(X_j)
 -H\!\left(\sum_{j=1}^m X_j\right)
 \leq 30m\sum_{j=1}^m\mathcal E(\operatorname{Law}(X_j)).
 }
\]
The authors identify this as an explicit quantitative strengthening of
the finite-law estimate underlying OpenAI's Lemma~3.2. In their
comparison, it implies a strengthened version of that lemma with
constant $C_m=120m^2$, whereas the original proof method yields a
constant of order $O(m^3)$. These constants refer to the corresponding
finite-law bounds and should not be confused with the $30m$ coefficient
in Theorem~1.3 itself.

\paragraph{Proof strategies and attribution.}
The arguments are different. OpenAI's Result~148 uses randomized
nested partitions, fair two-point submeasures, conditioning on block
symbol counts, and telescoping entropy gains at disjoint fine-scale
windows. Kittle--Kogler prove their entropy-loss inequality by Poisson
partitions and a Margulis--Russo differentiation argument, and deduce
the dimension formula using Hochman's entropy theorem and a
variance-summation method developed in earlier work of the authors.
Kittle--Kogler explicitly acknowledge OpenAI's independently announced
Result~148 and describe the finite-law estimates as the central new
ingredients of both proofs.

\paragraph{Chronology, machine assistance, and formalization.}
OpenAI dates Result~148 to September 24, 2026 and released it publicly
on October 6, 2026. Kittle--Kogler report discovering their initial
proof on September 27, 2026 and posted their paper on October 7, 2026.
Their machine-contribution statement credits GPT-6 Astra with the
initial proof and entropy inequality, and describes subsequent human
rewriting and an alternative proof of the dimension deduction.
They report an approximately 47,000-line Lean formalization, available
at \url{https://github.com/ckkogler/exact-overlaps-one-dim-lean}.
The OpenAI contribution separately links its own Lean scope and
comparator files. Neither dimension theorem, by itself, asserts
absolute continuity, and the existence of proof code should not be
mistaken for an independent review by this site.

\paragraph{Sources for this comparison.}
\begin{itemize}[leftmargin=*,nosep]
 \item Kittle--Kogler, manuscript (Theorems~1.1 and~1.3;
       Corollary~1.2; machine-contribution statement):
       \url{https://arxiv.org/pdf/2610.10511}.
 \item OpenAI, \emph{The entropy-rate dimension formula for
       self-similar measures on the line}, Result~148 (September 24, 2026):
       \url{https://github.com/openai/math/blob/main/preprints/The-entropy-rate-dimension-formula-for-self-similar-measures-on-the-line-September-24-2026/main.pdf}.
 \item Kittle--Kogler, Lean formalization:
       \url{https://github.com/ckkogler/exact-overlaps-one-dim-lean}.
\end{itemize}

\end{document}
